% CLAUDE REPORT::WORSTCASE-BRIDGE-0001::2026-11-08
\documentclass{article}
\usepackage{amsmath,amssymb,amsthm}
\newtheorem{prop}{Proposition}
\begin{document}
\title{Worst-case derivation cost $\overline W_R(N)$: audit of the bridge $\mathrm{SAT}\in\mathrm P\Rightarrow \overline W_R\le\mathrm{poly}$}
\date{2026-11-08. Status: P versus NP remains OPEN.}
\maketitle

\section*{Setup}
Let $R$ be a fixed derivation system with cost $\mathrm{cost}_R$, and
$\overline W_R(N)=\max_{F\in\mathrm{UNSAT},|F|\le N}\min_{\pi:F\vdash_R\bot}\mathrm{cost}_R(\pi)$.
The inner minimum is a property of $(R,F)$ alone; free choice of the optimal $\pi$ (no charge for search) is already built in. It is the same quantity as optimal proof size, so granting it changes nothing relative to classical proof complexity.
Call $\mathrm{Bridge}(R)$ the statement $\mathrm{SAT}\in\mathrm P\Rightarrow\overline W_R\le\mathrm{poly}$.

\section*{Results (all elementary; KNOWN THEOREM / PROVED)}
\begin{prop}[Cook--Reckhow, KNOWN]
If $R$ has polynomial-time checkable derivations, then $\overline W_R$ polynomial implies $\mathrm{NP}=\mathrm{coNP}$. Hence a superpolynomial lower bound for such $R$ proves $\mathrm{NP}\ne\mathrm{coNP}$ only if it is for \emph{every} such $R$; for one $R$ it proves nothing about P vs NP.
\end{prop}

\begin{prop}[Bridge is vacuous for systems with known lower bounds]
If $\overline W_R$ is known superpolynomial (resolution, cutting planes, bounded-depth Frege, \dots), then $\mathrm{Bridge}(R)\iff \mathrm{SAT}\notin\mathrm P$. So proving the bridge for such $R$ is equivalent to proving $\mathrm P\neq\mathrm{NP}$.
\end{prop}
\emph{Proof.} The conclusion is false, so the implication is equivalent to the negation of its hypothesis. \qed

\begin{prop}[Bridge for $R_M$ is trivial and circular]
For a decider $M$ for SAT running in time $T$, let $R_M$ accept the one-line proof ``$M(\neg F)$ rejects.'' Then $\overline W_{R_M}\le T$. But $R_M$ depends on $M$; a fixed $R$ with $\mathrm{Bridge}(R)$ for all $M$ must $p$-simulate all $R_M$ for poly-time $M$ correct on SAT.
\end{prop}

\begin{prop}[What the bridge actually needs]
$\mathrm{Bridge}(R)$ holds whenever $R$ is $p$-optimal (p-simulates every Cook--Reckhow system): under $\mathrm P=\mathrm{NP}$, $R_M$ is polynomially bounded and checkable (running $M$), hence so is $R$. Existence of a $p$-optimal propositional proof system is OPEN; it is known to follow from $\mathrm{E}=\mathrm{NE}$ (Krají\v cek--Pud\-l\'ak 1989) and is widely conjectured false. If $R$ is $p$-optimal, a superpolynomial $\overline W_R$ would give $\mathrm{NP}\ne\mathrm{coNP}$, hence $\mathrm P\ne\mathrm{NP}$.
\end{prop}

\section*{Consequence for the program}
\begin{enumerate}
\item The weakest sufficient bridge for a \emph{fixed} $R$ is, in effect, ``$R$ is conditionally $p$-optimal'': $\mathrm P=\mathrm{NP}\Rightarrow R\ p\text{-simulates every }R_M$. No known natural $R$ (Frege, extended Frege, SAT-circuit DAG systems) is known to satisfy it; extended Frege is the natural candidate, and its lower bounds are open.
\item The obstruction is the unverifiable-correctness gap: $M$'s correctness on all UNSAT inputs is not poly-time checkable inside $R$, unlike Levin search for satisfying assignments. Any design of a ``fixed combined inference system'' must internalize a poly-time simulation of arbitrary $M$ \emph{and} a poly-size certificate of $M$'s correctness; this is the reflection-principle formulation (Krajíček--Pudlák).
\item Hard-instance families surviving resolution and algebraic reasoning, even if found, give only $\overline W_R$ bounds for weak $R$ and do not meet the bridge by Proposition~2.
\end{enumerate}
Next target: decide whether the project's SAT-circuit DAG system $R^\ast$ $p$-simulates extended Frege (then it is at least as strong as the best natural candidate), and if so state Bridge$(R^\ast)$ as conditional on a stated reflection hypothesis.

\section*{Literature check}
Egidy, \emph{The SPARSE-Relativization Framework and Applications to Optimal Proof Systems}, arXiv:2602.02294 (Feb 2026): existence of $p$-optimal systems for TAUT still open. Krajíček--Pudlák 1989; Cook--Reckhow 1979. Transfer classification: Props.~1--4 are directly transferable theorems; the DAG-hypergraph state/work model is an analogy until its verifier is shown poly-time.

\textbf{Status: no lower bound or transfer theorem proved; P versus NP remains OPEN.}
% END_OF_VERIFIED_REPORT::CLAUDE::WORSTCASE-BRIDGE-0001
\end{document}
